% Rooting the Synoptic Tree: Unsupervised Direction Detection by Redactional Polarization
% Compile with: xelatex main.tex && bibtex main && xelatex main.tex && xelatex main.tex
\documentclass[11pt,letterpaper]{article}

% ── Core packages ───────────────────────────────────────────────────
\usepackage[utf8]{inputenc}
\usepackage{fontspec}
\usepackage{microtype}
\usepackage{booktabs}
\usepackage{graphicx}
\usepackage[table]{xcolor}
\usepackage{hyperref}
\usepackage{enumitem}
\usepackage{csquotes}
\usepackage[authoryear,round]{natbib}
\usepackage[letterpaper, top=0.9in, bottom=1.0in, left=1.0in, right=1.0in]{geometry}
\usepackage{titlesec}
\usepackage{fancyhdr}
\usepackage{tcolorbox}
\tcbuselibrary{breakable}
\usepackage{array}
\usepackage{tabularx}
\usepackage{ragged2e}
\usepackage{amsmath}
\usepackage{float}
\usepackage[font={small,color=stone},labelfont={bf,color=marine},labelsep=quad]{caption}

\setlist{nosep, leftmargin=*}
\frenchspacing
\setlength{\parindent}{0pt}
\setlength{\parskip}{0.55em}

% ── Typography ───────────────────────────────────────────────────────
\setmainfont{TeX Gyre Pagella}[Ligatures=TeX]
\newfontfamily\headfont{Poppins}
\newfontfamily\greekfont{FreeSerif}
\newcommand{\gk}[1]{{\greekfont\itshape #1}}
\setmonofont{Latin Modern Mono}[Scale=0.92]

% ── Palette ──────────────────────────────────────────────────────────
\definecolor{marine}{HTML}{123C48}
\definecolor{marinelight}{HTML}{2E6272}
\definecolor{terracotta}{HTML}{B75B36}
\definecolor{papyrus}{HTML}{F7F2E7}
\definecolor{stone}{HTML}{6E6A5F}
\definecolor{rulegrey}{HTML}{D9D2C0}
\definecolor{inkblack}{HTML}{201D1A}
\color{inkblack}
\arrayrulecolor{marine!65!black}

\newcolumntype{P}[1]{>{\centering\arraybackslash}p{#1}}
\newcolumntype{L}[1]{>{\RaggedRight\arraybackslash}p{#1}}

\renewcommand\labelitemi{\textcolor{terracotta}{\textbullet}}

\hypersetup{
  colorlinks=true,
  linkcolor=marine,
  citecolor=marine,
  urlcolor=terracotta,
  pdftitle={Rooting the Synoptic Tree: Unsupervised Direction Detection by Redactional Polarization},
  pdfauthor={Abderahmane Ainouche}
}

% ── Section styling ─────────────────────────────────────────────────
\titleformat{\section}
  {\headfont\Large\bfseries\color{marine}}
  {\color{marine}\thesection}
  {0.65em}{}
  [{\vspace{2pt}\color{rulegrey}\titlerule[0.9pt]}]
\titlespacing*{\section}{0pt}{1.7em}{0.9em}

\titleformat{\subsection}
  {\headfont\large\bfseries\color{marinelight}}
  {\thesubsection}{0.6em}{}
\titlespacing*{\subsection}{0pt}{1.3em}{0.5em}

% ── Header / footer ─────────────────────────────────────────────────
\pagestyle{fancy}
\fancyhf{}
\fancyhead[L]{\footnotesize\headfont\color{stone} Rooting the Synoptic Tree}
\fancyhead[R]{\footnotesize\headfont\color{stone} Redactional Polarization}
\fancyfoot[C]{\footnotesize\headfont\color{stone}\thepage}
\renewcommand{\headrulewidth}{0.4pt}
\renewcommand{\headrule}{{\color{rulegrey}\hrule height\headrulewidth}}
\fancypagestyle{firstpage}{%
  \fancyhf{}
  \fancyfoot[C]{\footnotesize\headfont\color{stone}\thepage}
  \renewcommand{\headrulewidth}{0pt}
}

% ── Small helpers ────────────────────────────────────────────────────
\newcommand{\statcard}[2]{%
  \cellcolor{papyrus}{\begin{minipage}[c][2.05cm][c]{\linewidth}\centering
    {\headfont\bfseries\fontsize{17}{19}\selectfont\color{marine}#1}\\[3pt]
    {\footnotesize\headfont\color{stone}#2}
  \end{minipage}}%
}
\newcommand{\tblhead}{\rowcolor{marine!12}}

\begin{document}

% ── Title block ──────────────────────────────────────────────────────
\thispagestyle{firstpage}
\begin{flushleft}
{\headfont\footnotesize\bfseries\color{terracotta}\MakeUppercase{Digital Humanities \textbullet\ Computational Textual Criticism}}\\[7pt]
{\headfont\bfseries\fontsize{21}{26}\selectfont\color{marine}%
Rooting the Synoptic Tree: Unsupervised Direction Detection\\ by Redactional Polarization}\\[10pt]
{\headfont\large\color{stone} Abderahmane Ainouche}
\end{flushleft}
\vspace{6pt}
{\color{rulegrey}\rule{\linewidth}{0.9pt}}
\vspace{1.3em}

% ── Abstract ─────────────────────────────────────────────────────────
\begin{tcolorbox}[
  breakable,
  colback=papyrus, colframe=marine,
  boxrule=0pt, leftrule=2.6pt, arc=0pt, outer arc=0pt,
  left=16pt, right=16pt, top=11pt, bottom=11pt
]
{\headfont\bfseries\color{marine} Abstract}\par\vspace{5pt}
\noindent The \textit{Synoptic Problem}---the question of how Matthew, Mark, and Luke are literarily related---turns on \textit{direction}: given two parallel passages, which is the source and which the copy? We show that this is precisely the stemmatological \textit{rooting problem}, and that \textit{global} passage-level scores (embedding similarity, conditional-generation codelength, and a learned redaction head) cannot solve it: each is confounded with passage length or with Mark's distinctive style, a fact we prove with two real known-direction corpora of opposite length polarity. We then introduce the \textbf{Redactional Polarization Model (RPM)}, which scores direction at the level of the individual \textit{variant} using the classical canons of textual criticism. A single length-free canon---\textit{connective smoothing} (the redactional replacement of paratactic \gk{καί} with a smoother connective)---reaches \textbf{0.78 directed accuracy} on the Synoptic triple tradition (0.88 on pericopes carrying such an edit; 0.97 on its most confident quartile via abstention), transfers to non-Synoptic known-direction data, and survives a strict control for Mark's global \gk{καί}-density. Pooling per-pericope votes into a Bayesian posterior over the four classical hypotheses \textbf{recovers Markan priority with no Synoptic supervision} (2SH~0.64, Farrer~0.36, Griesbach and Augustinian $\approx 0$), while the Farrer--Q question is left honestly open. Because RPM is unsupervised on the Synoptics, its posterior can enter a downstream Bayesian model comparison without circularity. All code and data are released under CC-BY-SA~4.0.
\end{tcolorbox}

\vspace{1.1em}
\begin{center}
\renewcommand{\arraystretch}{1}
\setlength{\tabcolsep}{7pt}
\begin{tabular}{P{0.18\linewidth}P{0.18\linewidth}P{0.18\linewidth}P{0.18\linewidth}P{0.18\linewidth}}
\statcard{0.78}{Directed acc.\ (synoptic)} &
\statcard{0.97}{@ 25\% coverage} &
\statcard{0.64}{2SH posterior} &
\statcard{$\approx$0}{Griesbach / Aug.} &
\statcard{3.5\,:\,1}{style-matched edits}
\end{tabular}
\end{center}
\vspace{0.6em}

\section{Introduction}
\label{sec:intro}

For more than two centuries, the \textit{Synoptic Problem} has asked how the first three
Gospels---Matthew, Mark, and Luke, which share extensive verbatim parallels---are literarily
related \citep{streeter1924}. The dominant \textbf{Two-Source Hypothesis} (2SH) holds that Mark
was written first and that Matthew and Luke each used Mark independently, together with a lost
sayings source \textbf{Q}. Its principal rival, the \textbf{Farrer hypothesis}, dispenses with Q
by having Luke use Matthew directly \citep{farrer1955,goodacre2001}; the older
\textbf{Griesbach} (Matthew first, Mark last) and \textbf{Augustinian} (Matthew first, Mark
second) hypotheses invert the priority entirely \citep{griesbach1978}.

What separates these hypotheses is not \textit{whether} the Gospels are related---they obviously
are---but the \textit{direction} of dependence: who copied whom. This is exactly the situation
that stemmatology, the study of manuscript genealogies, calls the \textbf{rooting problem}
\citep{trovato2017}. Recovering the unrooted tree (the pattern of relatedness) is easy; placing
the \textit{root} (assigning direction) is the hard part, and it is classically solved by
\textit{polarizing individual variants}---deciding, reading by reading, which is primitive and
which is derived---using the canons of textual criticism \citep{metzger2005}.

Prior computational work on the Synoptics has modelled patterns of word agreement
\citep{abakuks2006} but has not produced a direction detector validated against known-direction
data. In this paper we (i) show why the natural \textit{global} approaches fail, as a principled
negative result; (ii) introduce the \textbf{Redactional Polarization Model (RPM)}, a variant-level
detector grounded in the textual-criticism canons; and (iii) use it to root the Synoptic tree
without any Synoptic supervision. Because the detector never sees the disputed answer, its output
is admissible as evidence in a downstream Bayesian test of the hypotheses---a property we regard
as essential and that supervised classifiers trained on 2SH labels cannot claim.

\section{The Direction Problem and Its Built-in Trap}
\label{sec:trap}

Under any hypothesis, the labelled Synoptic data is confounded \textit{by construction}. In the
triple tradition---material present in all three Gospels---2SH makes Mark the source in every
pair, and treats every Matthew--Luke pair as \textit{independent} (both descend from Mark). Thus
``direction'' is collinear with ``is this Mark'': any classifier can score well by detecting
Mark's style rather than the direction of copying. A genuine detector must isolate direction
\textit{per se}, and must be validated where the answer is independently known and where a
Synoptic author is \textit{absent}.

We therefore anchor every claim on two external corpora of \textbf{known} direction and
\textbf{opposite length polarity}: \textbf{Jude~$\rightarrow$~2~Peter} (the later text
\textit{expands} its source) and the \textbf{Septuagint Samuel--Kings~$\rightarrow$~Chronicles}
synopsis (the later text \textit{abridges} its source). A real detector must give the right
answer on \textit{both}; a length or style artefact will not.

\section{Why Global Passage Scores Fail}
\label{sec:negatives}

We first tested the natural family of \textit{global} scores---each reducing a passage pair to a
single number---and found every one confounded. This is a contribution in its own right: it maps
the boundary of what passage-level information can and cannot reveal.

\begin{table}[H]
\centering
\small
\begin{tabular}{L{4.6cm}L{4.4cm}L{4.9cm}}
\toprule
\tblhead\textbf{Global score} & \textbf{Apparent signal} & \textbf{What it actually measures} \\
\midrule
Embedding cross-similarity (10 asymmetry features) & none (chance) & nothing usable \\
Conditional-generation codelength $\mathrm{NLL}(B\!\mid\!A)-\mathrm{NLL}(A\!\mid\!B)$ & strong on synoptics & Mark's style + length \\
Length-fair MDL / compression score & strong on synoptics, $\approx 0$ external & typicality + length \\
Learned redaction head (synthetic training) & 0.99 on synthetic authors & generator artefact + length prior \\
\bottomrule
\end{tabular}
\caption{\textbf{The global-score dead ends.} Every passage-level score reduces, under control,
to passage length or to Mark's (simpler, more predictable) Greek---not to copying direction.}
\label{tab:negatives}
\end{table}

The decisive control is the two-polarity test. The conditional-codelength asymmetry, which looks
like a strong direction signal on the Synoptics, \textit{reverses its sign} between
Jude~$\rightarrow$~2~Peter (copy longer) and Chronicles (copy shorter): it tracks \textit{which
text is longer}, not which is the source. The learned redaction head, which reaches 99\% on
held-out synthetic authors, collapses to a pure length prior on real data (17\% on the expanding
pair, 100\% on the abridging pair). The lesson is granularity: \textbf{averaging a passage into a
single number drowns the few directional edits under length and authorial style}. This motivates
scoring at the level of the individual variant.

\section{The Redactional Polarization Model}
\label{sec:rpm}

RPM treats direction as rooting by variant polarization. Given an aligned pair $(X,Y)$ it (i)
extracts typed \textbf{variants}---substitutions, and one-sided pluses---from the token
alignment; (ii) scores each variant with \textit{signed, local} features drawn from the
textual-criticism canons, with the convention that a positive value means \textit{$X$'s reading
is primitive} (i.e.\ $X$ is the source); and (iii) aggregates them into a per-pair score.

\subsection{Canon features}

Three canons yield a computable, signed, per-variant feature:
\begin{itemize}
\item \textbf{Lectio difficilior} (harder reading is older): the lexical-rarity (markedness)
  difference between $X$'s and $Y$'s readings, from a Koine frequency table.
\item \textbf{Lectio brevior}, applied \textit{locally} (the shorter reading is older): the
  length difference of \textit{this variant's} span---never the passage length.
\item \textbf{Connective smoothing}: a redactor tends to replace rough paratactic \gk{καί}
  (``and'') with a smoother connective (\gk{δέ}, \gk{οὖν}, \gk{τότε}, \gk{γάρ}); the rough side
  is scored primitive. This is a \textit{directional edit at an aligned position}, not a count.
\end{itemize}
Because each feature negates when $X$ and $Y$ are swapped, the aggregate is antisymmetric by
construction: swapping the inputs negates the score.

\subsection{Aggregation and abstention}

The per-pair score is a linear combination of summed variant features,
$S(X,Y)=\mathbf{w}^{\!\top}\!\sum_v \boldsymbol{\phi}(v)$, with $S>0$ read as ``$X$ is the
source''. Weights are fit by logistic regression on \textit{external} known-direction variants
(swap-augmented, no intercept, so antisymmetry is preserved); the sign of each weight is fixed a
priori by its canon. Confidence is $|S|$: on directionally neutral pericopes the score is near
zero and RPM \textbf{abstains} rather than guessing---matching how sparse the evidence truly is.
Crucially, no weight is ever fit on Synoptic data.

\section{Results}
\label{sec:results}

We report five gated findings, each with pericope-grouped (clustered) bootstrap 95\% confidence
intervals, so that the many pairs from one pericope are not treated as independent.

\subsection{Which canons polarize? (H1)}

Tested on the external known-direction variants across both length polarities: \textit{lectio
brevior} is a \textbf{pure length proxy} (0.00 accuracy on the abridging corpus, 1.00 on the
expanding one---caught red-handed) and is discarded. \textit{Lectio difficilior} is at chance.
\textbf{Connective smoothing is the survivor}: it is consistent on both polarities and reaches
0.80 on the Synoptics. All subsequent results use the length-free connective canon.

\subsection{Aggregation, abstention, and transfer (H2--H3)}

\begin{table}[H]
\centering
\begin{tabular}{lccc}
\toprule
\tblhead\textbf{Evaluation set} & \textbf{Directed acc.} & \textbf{@ 50\% cov.} & \textbf{@ 25\% cov.} \\
\midrule
\rowcolor{terracotta!12}\textbf{Synoptic} (Mark~$\rightarrow$~Matt/Luke) & \textbf{0.78 [0.71, 0.85]} & 0.93 & \textbf{0.97} \\
Jude~$\rightarrow$~2~Peter ($n=6$) & 1.00 & 1.00 & 1.00 \\
LXX Chronicles (in-domain) & 0.57 & --- & 0.64 \\
\bottomrule
\end{tabular}
\caption{\textbf{RPM directed accuracy with abstention.} Restricting to the most confident
pericopes sharply raises accuracy---the confident quartile is 97\% correctly directed. Dropping
the length feature entirely \textit{kept} the Synoptic result (0.76~$\rightarrow$~0.78), proving
the signal is connective smoothing, not length. Connective smoothing is a Gospel-genre
phenomenon, hence weak on Septuagintal narrative.}
\label{tab:rpm}
\end{table}

On the 121 Synoptic pericopes that actually carry a connective edit, directed accuracy is
\textbf{0.876}; the headline 0.78 is diluted only by the pericopes with no such edit, on which
RPM correctly abstains (they score zero and are dropped first, yielding 0.93 and 0.97 at 50\% and
25\% coverage).

\subsection{Is it direction, or Mark's \gk{καί}-heavy style? (control)}

The single most serious objection is that RPM votes ``Mark is the source'' merely because Mark
uses \gk{καί} more \textit{globally}. We control for this directly: stratifying the Mark--$X$
pericopes by $|\Delta|$, the gap in global \gk{καί}-density between the two texts. If the signal
were pure style, the \textit{matched} stratum would collapse to chance.

\begin{table}[H]
\centering
\small
\begin{tabular}{lccc}
\toprule
\tblhead\textbf{Stratum (by $|\Delta\,$\gk{καί}$|$)} & \textbf{Markan-priority acc.} & \textbf{95\% CI} & \textbf{$n$} \\
\midrule
Overall & 0.876 & [0.82, 0.93] & 121 \\
\rowcolor{papyrus}\textbf{$T_1$ --- matched \gk{καί}-density} & \textbf{0.784} & \textbf{[0.64, 0.90]} & 37 \\
$T_2$ --- middle & 0.950 & [0.87, 1.00] & 40 \\
$T_3$ --- divergent style & 0.886 & [0.78, 0.98] & 44 \\
\bottomrule
\end{tabular}
\caption{\textbf{The style-confound control.} Even where Mark and the other Gospel have matched
global \gk{καί}-density, the signal stays at 0.78 (CI clears chance). The per-edit vote ratio
(Mark-source\,:\,other-source) is 5.6\,:\,1 across all edits and 3.5\,:\,1 on the matched half.
The Markan-priority signal is therefore genuine \textit{per-edit} direction (\gk{καί} smoothed to
\gk{δέ}/\gk{τότε} at aligned positions), not global style.}
\label{tab:control}
\end{table}

\subsection{Does editorial fatigue add a second canon? (H4)}

Editorial fatigue---a copyist who makes characteristic changes early, then tires and reverts to
the source---is directional by nature \citep{goodacre1998}. We tested its one length-robust
feature (shared entities introduced later in the copy) as a second RPM canon. It does \textit{not}
help on the Synoptics: fitting a blend on the external corpus over-weights fatigue (the stronger
in-sample canon there) and \textit{collapses} the Synoptic result to 0.41, because aggregate
fatigue is at chance on the Synoptics (0.44, CI includes 0.5). A decisive test confirms it: on the
pericopes where the connective canon is \textit{silent}, fatigue is still at chance (0.44), so it
adds no complementary coverage. Aggregate fatigue is a real but \textit{genre-limited} canon
(useful on Septuagintal abridgement), complementary to connective smoothing rather than additive
on the Gospels; capturing it on the Synoptics would require a sharper, sparse dangling-reference
operator rather than an average.

\subsection{Rooting the tree, and the Farrer--Q question (H5)}

The four hypotheses are four rootings, each predicting a direction for the three pairwise
relationships (or, for 2SH on Matthew--Luke, \textit{independence}). Pooling the unsupervised
connective vote per pericope into a Beta--Bernoulli marginal likelihood per relationship---with
the 2SH ``independent'' prediction modelled explicitly as $\theta=0.5$---yields a posterior over
the four stemmata and a Bayes factor for Farrer versus Q \citep{kass1995}.

\begin{table}[H]
\centering
\small
\begin{minipage}[t]{0.54\linewidth}
\centering
\begin{tabular}{lccc}
\toprule
\tblhead\textbf{Relationship} & \textbf{Source vote} & \textbf{$n$} \\
\midrule
Matthew--Mark (triple) & Mark 75\% & 59 \\
Mark--Luke (triple) & Mark 91\% & 57 \\
Matthew--Luke (double / Q) & Matt 58\% & 12 \\
\bottomrule
\end{tabular}
\end{minipage}\hfill
\begin{minipage}[t]{0.42\linewidth}
\centering
\begin{tabular}{lc}
\toprule
\tblhead\textbf{Hypothesis} & \textbf{Posterior} \\
\midrule
\rowcolor{terracotta!12}\textbf{2SH} & \textbf{0.64} \\
Farrer & 0.36 \\
Griesbach & $\approx 0.00$ \\
Augustinian & $\approx 0.00$ \\
\bottomrule
\end{tabular}
\end{minipage}
\caption{\textbf{Unsupervised rooting.} The canon sees Mark as the source on both Markan
relationships (75\%, 91\%), so the Matthean-priority hypotheses (Griesbach, Augustinian)---which
require Mark to be a \textit{copy}---are excluded. 2SH and Farrer agree on Markan priority and
differ only on Matthew--Luke.}
\label{tab:rooting}
\end{table}

Two conclusions follow. First, \textbf{Markan priority is recovered with zero Synoptic
supervision}: the field-consensus result emerges from a textual-criticism canon validated only on
external corpora. Second, \textbf{Farrer versus Q is left open, leaning weakly to Q}. On the
double tradition---the Q material, where no Markan parallel exists---the canon finds no consistent
Matthew~$\rightarrow$~Luke direction (7 of 12 $\approx$ chance), which is what independence
predicts; the Bayes factor Farrer\,:\,2SH is 0.56 (anecdotal, favouring Q). The strong
Matthew-source signal on \textit{triple-tradition} Matthew--Luke ($\mathrm{BF}\approx165$) is not
valid Farrer evidence: there both evangelists used Mark, so a Matthew-source polarity is expected
under \textit{every} hypothesis. The clean test is the double tradition, and it is simply too
sparse in connective edits (12 pericopes) to settle the question.

\section{Discussion}
\label{sec:discussion}

\subsection{What has been shown}

RPM is, to our knowledge, the first computational method to recover Synoptic copying direction
from the text itself and to validate it against known-direction data of both length polarities.
It is interpretable---its weights are the textual-criticism canons---length-controlled, and
abstention-calibrated, and it reproduces Markan priority while excluding the Matthean-priority
hypotheses. Equally important is what it does \textit{not} claim: the Farrer--Q question, which
turns on a small and edit-sparse body of double-tradition material, remains open, and we report it
as such rather than overreading a near-chance result.

\subsection{Non-circularity and the road to a Bayesian verdict}

Because RPM is unsupervised on the Synoptics---its canon signs are fixed a priori and its weights
are fit only on external corpora---its per-pericope log-odds may be pooled into a Bayesian model
comparison over the four hypotheses \textit{without} the circularity that would invalidate a
classifier trained on 2SH labels. The rooting posterior of Table~\ref{tab:rooting} is a first,
deliberately simple instance of that pooling; a full hierarchical treatment with prior-sensitivity
analysis is the natural next step.

\subsection{Limitations}

Three limits bound the claims. (i) Connective smoothing is a Gospel-genre phenomenon; it transfers
only partially to Septuagintal narrative, and a richer canon set is needed for cross-genre
generality. (ii) The Farrer--Q test is underpowered; resolving it needs more double-tradition
data or a sharper, sparse fatigue operator than the aggregate tested here. (iii) While the
per-edit control (Table~\ref{tab:control}) shows the Markan-priority signal is not reducible to
Mark's global \gk{καί}-density, connective usage remains a stylistic as well as a redactional
variable, and independent canons agreeing on the same root would further harden the conclusion.

\section{Conclusion}
\label{sec:conclusion}

We reframed Synoptic direction detection as the stemmatological rooting problem and showed that
global passage scores cannot solve it, while variant-level polarization by the textual-criticism
canons can. The Redactional Polarization Model reaches 0.78 directed accuracy on the Synoptic
triple tradition (0.97 on its confident quartile), survives a strict control for Mark's style, and
---pooled over the corpus---recovers Markan priority and excludes the Matthean-priority hypotheses
with no Synoptic supervision, leaving the Farrer--Q question honestly open. Code and data are
released under CC-BY-SA~4.0 to support further work in computational textual criticism.

% ── References ──────────────────────────────────────────────────────
\bibliographystyle{apalike}
\bibliography{references}

\end{document}
